\documentclass[12pt,a4paper]{article}
\usepackage[utf8]{inputenc}
\usepackage[T1]{fontenc}
\usepackage[vietnamese]{babel} % Hỗ trợ hiển thị tiếng Việt chuẩn
\usepackage{geometry}
\usepackage{graphicx}
\usepackage{hyperref}
\usepackage{listings}
\usepackage{xcolor}
\usepackage{booktabs}
\usepackage{tabularx}
\usepackage{enumitem}
\usepackage{fancyhdr}
\usepackage{titlesec}
\usepackage{tcolorbox}
\usepackage{float}
\usepackage{multirow}
\usepackage{amsmath}
\usepackage{setspace} % Cần cho lệnh \setstretch
\geometry{top=2.5cm, bottom=2.5cm, left=3cm, right=2.5cm}
\definecolor{primary}{RGB}{30,60,120}
\definecolor{secondary}{RGB}{200,169,110}
\definecolor{codeblue}{RGB}{0,112,192}
\definecolor{codegray}{RGB}{128,128,128}
\definecolor{codegreen}{RGB}{0,128,0}
\definecolor{codebg}{RGB}{245,245,245}
\definecolor{lightblue}{RGB}{219,234,254}
\lstset{
backgroundcolor=\color{codebg},
commentstyle=\color{codegreen},
keywordstyle=\color{codeblue}\bfseries,
stringstyle=\color{orange},
basicstyle=\ttfamily\footnotesize,
breaklines=true,
frame=single,
rulecolor=\color{gray!30},
numbers=left,
numberstyle=\tiny\color{codegray},
numbersep=5pt,
tabsize=2
}
\pagestyle{fancy}
\fancyhf{}
\fancyhead[L]{\small\textcolor{gray}{Graph RAG Chatbot -- Quy chế đào tạo ĐHBK Hà Nội}}
\fancyhead[R]{\small\textcolor{gray}{\thepage}}
\fancyfoot[C]{\small\textcolor{gray}{Báo cáo Project 2}}
\renewcommand{\headrulewidth}{0.4pt}
\renewcommand{\footrulewidth}{0.4pt}
\titleformat{\section}{\Large\bfseries\color{primary}}{\thesection.}{0.5em}{}[\titlerule]
\titleformat{\subsection}{\large\bfseries\color{primary!80}}{\thesubsection.}{0.5em}{}
\titleformat{\subsubsection}{\normalsize\bfseries\color{primary!60}}{\thesubsubsection.}{0.5em}{} 
\hypersetup{colorlinks=true, linkcolor=primary, urlcolor=codeblue}

\begin{document}
%% ============================================================
%% TRANG BÌA
%% ============================================================
\begin{titlepage}
\centering
\setstretch{1.2}

{\Large\bfseries ĐẠI HỌC BÁCH KHOA HÀ NỘI} \\
\vspace{0.2cm}
{\Large\bfseries TRƯỜNG CÔNG NGHỆ THÔNG TIN VÀ TRUYỀN THÔNG} \\
\vspace{1.5cm}

{\Large\bfseries \rule{8cm}{0.4pt}} \\[0.5cm]
{\Huge\bfseries \color{primary} Báo cáo bài tập lớn} \\[0.5cm]

\vspace{1.5cm}

% Phần Logo - Nhớ upload file logobk.png lên Overleaf
\begin{figure}[h]
    \centering
    \includegraphics[width=4cm]{logobk.png}
\end{figure}

\vspace{1.5cm}

{\Large\bfseries Học phần: Project 2} \\[0.3cm]
{\large\bfseries Đề tài: Ứng dụng Graph RAG trong xây dựng chatbot hỏi đáp quy chế} \\[0.9cm]
\vspace{1.5cm}
{\large\bfseries Họ và tên: Bùi Công Hoan} \\[0.3cm]
{\large\bfseries MSSV: 20235083}  \\
\end{titlepage}

\tableofcontents
\newpage 

%% ============================================================
\section{Giới thiệu dự án}
\subsection{Bối cảnh và Động lực}
Quy chế đào tạo của Đại học Bách Khoa Hà Nội (QCDT\_2025) là văn bản pháp lý quan trọng quy định toàn bộ hoạt động học tập của sinh viên và nghiên cứu sinh. 
Với 48 điều khoản trải dài qua 6 chương, tài liệu này chứa đựng hàng trăm quy định chi tiết về điểm số, tín chỉ, điều kiện tốt nghiệp, cảnh báo học tập và nhiều vấn đề quan trọng khác. Thực tế cho thấy nhiều sinh viên gặp khó khăn khi tra cứu thông tin đó khi cần tìm hiểu thông tin hoặc tham gia làm bài kiểm tra quy chế:
\begin{itemize}
\item Văn bản dài, nhiều điều khoản liên kết chéo nhau.
\item Ngôn ngữ pháp lý khó hiểu với sinh viên, đặc biệt là với sinh viên năm nhất.
\item Mất nhiều thời gian để tìm kiếm thông tin cụ thể.
\item Các chatbot thông thường không đủ chính xác cho văn bản pháp quy và hay gặp tình trạng "ảo giác" (hallucination).
\end{itemize}

\subsection{Mục tiêu dự án}
Dự án xây dựng hệ thống Graph RAG Chatbot - một chatbot thông minh có khả năng:
\begin{itemize}
\item Trả lời chính xác các câu hỏi về quy chế đào tạo.
\item Trích dẫn đúng điều khoản từ văn bản gốc, cung cấp thông số/chỉ số chuẩn xác.
\item Hiểu được mối quan hệ phức tạp giữa các quy định.
\item Hoạt động 100\% nội bộ (offline), bảo mật dữ liệu tuyệt đối.
\end{itemize}

\subsection{Điểm nổi bật: Graph RAG và RAG truyền thống}
Khác với RAG (Retrieval-Augmented Generation) thông thường chỉ tìm kiếm theo độ tương đồng vector (vector similarity), Graph RAG kết hợp thêm duyệt đồ thị tri thức (graph traversal) - cho phép tìm được các thông tin liên quan gián tiếp mà vector search thường xuyên bỏ sót.

Ví dụ: Khi hỏi về "điều kiện học cùng lúc 2 chương trình", hệ thống không chỉ tìm node đó mà còn duyệt sang các node vệ tinh lân cận như "CPA", "Xếp hạng năm học", "Điều kiện tốt nghiệp" để bổ sung ngữ cảnh trọn vẹn cho câu trả lời.

%% ============================================================
\section{Kiến trúc hệ thống}
\subsection{Tổng quan kiến trúc}
Hệ thống được thiết kế theo kiến trúc 3 tầng rõ ràng:
\begin{table}[H]
\centering
\caption{Các tầng kiến trúc hệ thống}
\begin{tabularx}{\textwidth}{lXl}
\toprule
\textbf{Tầng} & \textbf{Thành phần} & \textbf{Công nghệ} \\
\midrule
Presentation & Giao diện chat, Graph visualization & Next.js 16.2, TypeScript, D3.js \\
Application & REST API, WebSocket, Hybrid Retriever, Generator & FastAPI, Python \\
Data & Graph database, Vector index & Neo4j 5.20, SBERT \\
\bottomrule
\end{tabularx}
\end{table}

\subsection{Các thành phần chính}
\subsubsection{Ingestion Pipeline (Xử lý dữ liệu đầu vào)}
\begin{itemize}
\item \textbf{PDF Loader}: Đọc file quy chế, trích xuất văn bản (chunking theo từng Điều).
\item \textbf{Entity Extractor}: Dùng LLM (Gemini) bóc tách thực thể (entities) và các mối quan hệ (relations).
\item \textbf{Embedder}: Tạo vector embeddings 768 chiều bằng mô hình \texttt{keepitreal/vietnamese-sbert}.
\item \textbf{Neo4j Loader}: Nạp nodes, edges và thiết lập vector index vào cơ sở dữ liệu Neo4j.
\end{itemize}

\subsubsection{Hybrid Retriever (Tìm kiếm lai)}
\begin{itemize}
\item \textbf{Vector Search}: Tìm 10 "Seed nodes" gần nghĩa nhất với câu hỏi bằng cosine similarity.
\item \textbf{Graph Traversal (1-hop)}: Sử dụng ngôn ngữ Cypher để duyệt sang các node hàng xóm trực tiếp của Seed nodes. Xếp hạng ưu tiên theo mật độ liên kết (\texttt{ORDER BY connection\_count DESC}) và trích xuất tối đa 25 nodes vệ tinh.
\item \textbf{Context Formatting}: Lọc và ghép nội dung của toàn bộ Nodes + Relations thành ngữ cảnh (Context) hoàn chỉnh gửi cho LLM.
\end{itemize}

\subsubsection{Generator (Sinh văn bản)}
Sử dụng model \texttt{qwen2.5:7b} chạy local qua Ollama:
\begin{itemize}
\item Streaming token-by-token mượt mà qua WebSocket.
\item Session memory (lưu 20 turns gần nhất để hiểu ngữ cảnh trò chuyện).
\item System Prompt nghiêm ngặt: Yêu cầu trả lời 100\% Tiếng Việt, cấm rào đón, cấm bịa đặt và bắt buộc phải rút trích chính xác các con số (tín chỉ, điểm, tháng) có trong ngữ cảnh.
\end{itemize}

\subsection{Luồng xử lý một câu hỏi}
\[
\text{Câu hỏi} \xrightarrow{\text{SBERT}} \text{Vector 768D} \xrightarrow{\text{Neo4j}} \text{10 Seeds}
\xrightarrow{\text{1-hop Graph}} \text{35 Nodes} \xrightarrow{\text{LLM (Qwen)}} \text{Câu trả lời}
\]

%% ============================================================
\section{Knowledge Graph}
\subsection{Thiết kế đồ thị tri thức}
Knowledge Graph là thành phần cốt lõi phân biệt hệ thống chatbot này với chatbot thông thường khác.
\begin{table}[H]
\centering
\caption{Thống kê Knowledge Graph}
\begin{tabular}{lr}
\toprule
\textbf{Chỉ số} & \textbf{Giá trị} \\
\midrule
Nodes & 286 \\
Edges & 374 \\
Node types  & 5 \\
Relation types  & 45 \\
Vector dimensions & 768 \\
\bottomrule
\end{tabular}
\end{table}

\subsection{Phân loại Nodes}
\begin{table}[H]
\centering
\caption{Phân loại nodes trong Knowledge Graph}
\begin{tabularx}{\textwidth}{llX}
\toprule
\textbf{Loại} & \textbf{Ví dụ} & \textbf{Mô tả} \\
\midrule
\texttt{Dieu} & Phạm vi điều chỉnh & Các điều khoản tổng quan \\
\texttt{KhaiNiem} & Tín chỉ, GPA, CPA & Khái niệm, định nghĩa trong quy chế \\
\texttt{DieuKien} & Điều kiện tốt nghiệp & Yêu cầu, điều kiện cụ thể \\
\texttt{QuyTrinh} & Đăng ký học tập & Quy trình, thủ tục\\
\texttt{DoiTuong} & Sinh viên, NCS & Đối tượng áp dụng \\
\bottomrule
\end{tabularx}
\end{table}

\subsection{Các quan hệ chính (Edges)}
\begin{table}[H]
\centering
\caption{Các loại quan hệ phổ biến nhất}
\begin{tabular}{llr}
\toprule
\textbf{Quan hệ} & \textbf{Ý nghĩa} & \textbf{Số lượng} \\
\midrule
\texttt{LIEN\_QUAN} & A liên quan đến B & 65 \\
\texttt{THUOC} & A thuộc về/là thành phần của B & 57 \\
\texttt{DAN\_DEN} & A dẫn đến hệ quả B & 39 \\
\texttt{YEU\_CAU} & A yêu cầu phải có B & 30 \\
\texttt{ANH\_HUONG} & A ảnh hưởng đến B & 16 \\
\texttt{AP\_DUNG\_CHO} & A áp dụng cho đối tượng B & 15 \\
\texttt{QUYET\_DINH} & A quyết định/xác định B & 15 \\
\bottomrule
\end{tabular}
\end{table}

%% ============================================================
\section{Công nghệ sử dụng}
\subsection{Backend}
\begin{table}[H]
\centering
\caption{Công nghệ Backend}
\begin{tabularx}{\textwidth}{llX}
\toprule
\textbf{Công nghệ} & \textbf{Phiên bản} & \textbf{Vai trò} \\
\midrule
Python & 3.11.9 & Ngôn ngữ lập trình chính \\
FastAPI & 0.111.0 & Web framework, REST API + WebSocket \\
Neo4j Driver & 5.20.0 & Kết nối Neo4j, thực thi Cypher queries \\
Ollama & 0.6.2 & Chạy LLM local (qwen2.5:7b) \\
sentence-transformers & 3.0.0 & Vietnamese-SBERT embedding \\
pypdf & 4.2.0 & Đọc và trích xuất văn bản PDF \\
pydantic-settings & 2.3.0 & Quản lý cấu hình, validation \\
uvicorn & 0.30.0 & ASGI server cho FastAPI \\
\bottomrule
\end{tabularx}
\end{table}

\subsection{Frontend}
\begin{table}[H]
\centering
\caption{Công nghệ Frontend}
\begin{tabularx}{\textwidth}{llX}
\toprule
\textbf{Công nghệ} & \textbf{Phiên bản} & \textbf{Vai trò} \\
\midrule
Next.js & 16.2.4 & React framework, App Router \\
TypeScript & 5.x & Type-safe JavaScript \\
Tailwind CSS & 4.x & Utility-first CSS framework \\
D3.js & 7.9.0 & Graph visualization, tương tác đồ thị \\
lucide-react & 1.11.0 & Icon library \\
\bottomrule
\end{tabularx}
\end{table}

\subsection{Hạ tầng (Infrastructure)}
\begin{table}[H]
\centering
\caption{Hạ tầng và công cụ}
\begin{tabularx}{\textwidth}{llX}
\toprule
\textbf{Công nghệ} & \textbf{Phiên bản} & \textbf{Vai trò} \\
\midrule
Neo4j & 5.20.0 & Graph database, vector index cosine similarity \\
Docker Desktop & 4.x & Container hóa Neo4j \\
Docker Compose & v5.1.2 & Orchestration các service \\
Ollama & latest & LLM runtime local, quản lý model \\
Node.js & v22.17.1 & Runtime cho Next.js \\
\bottomrule
\end{tabularx}
\end{table}

\subsection{AI Models}
\subsubsection{qwen2.5:7b (via Ollama)}
\begin{itemize}
\item \textbf{Nhà phát triển}: Alibaba Cloud (Qwen Team).
\item \textbf{Kích thước}: File weights $\sim$4.7 GB.
\item \textbf{Điểm mạnh}: Hiểu tiếng Việt cực tốt, reasoning chính xác, tuân thủ chặt chẽ System Prompt.
\item \textbf{Khắc phục lỗi}: Đã tùy chỉnh Prompt để khóa ngôn ngữ (ép tiếng Việt 100\%) nhằm tránh tình trạng mô hình rơi vào "ảo giác" trả lời bằng tiếng Trung.
\end{itemize}

\subsubsection{keepitreal/vietnamese-sbert}
\begin{itemize}
\item \textbf{Loại}: Sentence-BERT fine-tuned chuyên biệt cho tiếng Việt.
\item \textbf{Đầu ra}: Vector 768 chiều.
\item \textbf{Kích thước}: $\sim$540 MB.
\item \textbf{Sử dụng}: Tạo embeddings cho nodes đồ thị và biểu diễn câu hỏi người dùng thành không gian vector.
\end{itemize}

%% ============================================================
\section{Quy trình thực hiện}
\subsection{Giai đoạn 1: Setup môi trường}
Các bước đã thực hiện:
\begin{itemize}
\item Cài đặt Python 3.11.9, Node.js v22.17.1, Docker Desktop 4.x, Ollama.
\item Tạo cấu trúc thư mục project theo kiến trúc thiết kế.
\item Tạo Python virtual environment, cài đặt toàn bộ dependencies.
\item Cấu hình file \texttt{.env} với các biến môi trường (Neo4j, AI model).
\item Khởi động Neo4j 5.20.0 qua Docker Compose.
\item Khởi tạo Next.js 16 với TypeScript, Tailwind CSS, tích hợp D3.js.
\item Xác nhận kết nối với Neo4j Browser tại \texttt{http://localhost:7474}.
\end{itemize}

\subsection{Giai đoạn 2: Ingestion Pipeline}
Dữ liệu được tạo ra thông qua Entity Extraction Pipeline kết hợp Gemini (ở khâu nháp) và được tinh chỉnh, soát lỗi thủ công toàn bộ 31 trang PDF quy chế để đưa vào file \texttt{entities\_cache.json}. Phương pháp này đảm bảo:
\begin{itemize}
\item \textbf{Độ chính xác tuyệt đối}: 100\% thông tin được kiểm duyệt khớp với văn bản gốc.
\item \textbf{Bao phủ đầy đủ}: 48/48 điều khoản, mỗi điều tối thiểu 3 nodes vệ tinh.
\item \textbf{Quan hệ có ý nghĩa}: Các edges phản ánh đúng logic pháp lý (Ví dụ: "Học phần A -> Học trước -> Học phần B").
\item \textbf{Bảo toàn dữ liệu số}: Các ngưỡng (số TC, điểm, thời hạn) được trích xuất nguyên vẹn phục vụ trả lời.
\end{itemize}

\subsection{Giai đoạn 3: Backend FastAPI}
\subsubsection{API Endpoints}
\begin{table}[H]
\centering
\caption{Danh sách API Endpoints}
\begin{tabular}{lll}
\toprule
\textbf{Method} & \textbf{Endpoint} & \textbf{Chức năng} \\
\midrule
POST & \texttt{/chat} & Trả lời Chat (REST API non-streaming) \\
WebSocket & \texttt{/ws/chat} & Trả lời Chat streaming token-by-token \\
GET & \texttt{/graph/nodes} & Lấy toàn bộ danh sách Nodes \\
GET & \texttt{/graph/edges} & Lấy toàn bộ danh sách Edges \\
GET & \texttt{/graph/stats} & Thống kê đồ thị (tổng số node, edge) \\
GET & \texttt{/health} & Health check hệ thống \\
\bottomrule
\end{tabular}
\end{table}

\subsection{Giai đoạn 4: Frontend Next.js}
\begin{tcolorbox}[colback=green!5, colframe=codegreen, title=\textbf{Trạng thái}]
Hoàn thành -- Giao diện mượt mà, tích hợp đầy đủ tính năng Graph Visualization.
\end{tcolorbox}
Các component đã xây dựng:
\begin{itemize}
\item \textbf{Header}: Logo, tên ứng dụng, badge trạng thái Online.
\item \textbf{MessageBubble}: Hiển thị tin nhắn người dùng và bot, hỗ trợ streaming mượt mà.
\item \textbf{ChatInput}: Box nhập liệu, nút gửi, block thao tác khi đang chờ.
\item \textbf{GraphPanel}: Ứng dụng D3.js mô phỏng trực quan các Nodes/Edges mà bot đang suy luận để trả lời câu hỏi hiện tại.
\end{itemize}

%% ============================================================
\section{Kết quả và Đánh giá (Evaluation)}

\subsection{Hệ thống chấm điểm tự động (Evaluate Script)}
Để đo lường hiệu quả một cách khách quan, dự án đã xây dựng một bộ Script đánh giá dựa trên 50 câu hỏi test bao trùm nhiều thể loại (Điểm số, Đăng ký học, Điều kiện tốt nghiệp, Hallucination...).

\begin{table}[H]
\centering
\caption{Kết quả chạy Evaluation Local (Qwen 7B)}
\begin{tabularx}{\textwidth}{lX}
\toprule
\textbf{Chỉ số (Metric)} & \textbf{Kết quả đạt được} \\
\midrule
\textbf{Pass / Fail Ratio} & 43 Pass / 7 Fail (\textbf{86.0\% Pass Rate}) \\
\textbf{Node Hit Rate (Graph)} & \textbf{81.7\%} (Tìm trúng 81.7\% node chứa số liệu trọng tâm) \\
\textbf{Điều khoản Recall} & \textbf{100\%} (47/47 câu trích dẫn đúng điều khoản quy chế) \\
\textbf{Semantic Similarity} & 71.8\% (So với Ground Truth) \\
\textbf{Keyword Match} & 61.2\% (Lấy đúng các con số khó) \\
\textbf{Hallucination Control} & \textbf{100\% Pass} (3/3 câu hỏi bẫy bị hệ thống từ chối an toàn) \\
\textbf{Tốc độ trung bình} & $\sim$18.7 giây/câu (Tốc độ tối ưu cho Local Model chạy thuần CPU) \\
\bottomrule
\end{tabularx}
\end{table}

\subsection{Tổng kết Trạng thái Dự án}
\begin{table}[H]
\centering
\caption{Tổng kết các hạng mục công việc}
\begin{tabularx}{\textwidth}{lXc}
\toprule
\textbf{Hạng mục} & \textbf{Chi tiết} & \textbf{Trạng thái} \\
\midrule
Setup môi trường & Docker, Neo4j, Python venv, Next.js & \textcolor{codegreen}{Xong} \\
Knowledge Graph & 286 nodes, 374 edges, 48/48 điều & \textcolor{codegreen}{Xong} \\
Vector Search & Vietnamese-SBERT 768 chiều & \textcolor{codegreen}{Xong} \\
Graph Traversal & 1-hop traversal, ranking by connections & \textcolor{codegreen}{Xong} \\
Backend API & FastAPI + WebSocket streaming & \textcolor{codegreen}{Xong} \\
LLM Integration & Ollama qwen2.5:7b local (Ép System Prompt nghiêm ngặt) & \textcolor{codegreen}{Xong} \\
Frontend UI & Next.js dark theme chatbot & \textcolor{codegreen}{Xong} \\
Graph Visualization & D3.js interactive & \textcolor{codegreen}{Xong} \\
Production Deploy & Docker + Nginx + SSL & \textcolor{red}{Chưa} \\
\bottomrule
\end{tabularx}
\end{table}

\subsection{Ưu điểm}
\begin{enumerate}
\item \textbf{Miễn phí hoàn toàn}: Chạy 100\% local (offline), không phụ thuộc bất cứ API trả phí nào.
\item \textbf{Độ chính xác vượt trội}: Pass Rate lên đến 86\% với model 7B. Node Hit Rate cải thiện triệt để nhờ tối ưu Cypher Query.
\item \textbf{Trích dẫn minh bạch}: Luôn trích dẫn chính xác Số/Tên Điều khoản.
\item \textbf{Hiểu quan hệ logic}: Graph traversal kéo được các thông số phụ trợ (như con số, tháng, điểm liệt) xung quanh vấn đề chính.
\item \textbf{Trực quan hóa đồ thị}: Cửa sổ D3.js cho phép xem trực tiếp Bot đang đọc các Node nào trong não bộ để trả lời (Explainable AI).
\item \textbf{Kiểm soát ảo giác}: Không bịa chuyện (0\% Hallucination), trả lời nghiêm túc bằng 100\% tiếng Việt.
\end{enumerate}

\subsection{Hạn chế}
\begin{enumerate}
\item \textbf{Knowledge Graph khởi tạo thủ công}: Tốn thời gian rà soát, khó tự động cập nhật ngay khi quy chế thay đổi.
\item \textbf{Giới hạn Model nhỏ}: Qwen2.5:7b đôi khi chưa đủ năng lực logic (reasoning) để xử lý các câu hỏi lồng ghép nhiều tình huống giả định phức tạp.
\item \textbf{Phụ thuộc RAM}: Cần tối thiểu 8GB RAM (khuyến nghị 16GB) để duy trì cùng lúc Neo4j, SBERT, FastAPI và Ollama.
\item \textbf{Chưa có Production Deployment}: Hiện tại chỉ chạy tối ưu trên môi trường cục bộ (Local).
\end{enumerate}

\subsection{Hướng Phát triển}
\begin{enumerate}
\item \textbf{Nâng cấp Model}: Triển khai các Model Local lớn hơn (14B, 32B) hoặc kết nối API (OpenAI, Gemini) nếu dự án có kinh phí để nâng Pass Rate lên gần tuyệt đối (100\%).
\item \textbf{Dynamic Ingestion}: Phát triển module cho phép người dùng upload thêm file (Quy chế học bổng, KTX) và tự động nối Graph mới vào Graph cũ.
\item \textbf{Deploy Production}: Đóng gói toàn bộ hệ thống bằng Docker Swarm hoặc Kubernetes, triển khai trên VPS có GPU.
\end{enumerate}

%% ============================================================
\section{Hướng dẫn Cài đặt và Chạy}
\subsection{Yêu cầu Hệ thống}
\begin{itemize}
\item OS: Windows 10/11 (64-bit) hoặc Linux/macOS.
\item RAM: Tối thiểu 8GB, khuyến nghị 16GB.
\item Disk: Tối thiểu 15GB trống.
\item Cài đặt sẵn: Python 3.11+, Node.js 18+, Docker Desktop, Ollama.
\end{itemize}

\subsection{Khởi động Hệ thống}
\begin{lstlisting}[language=bash, caption=Các bước khởi động hệ thống]
# Bước 1: Khởi động Neo4j (Docker)
cd graphrag-quyche
docker compose up -d neo4j

# Bước 2: Kiểm tra Ollama
ollama list   # Đảm bảo đã pull qwen2.5:7b

# Bước 3: Chạy Backend
cd backend
.\venv\Scripts\activate
python -m uvicorn app.main:app --port 8000 --reload

# Bước 4: Chạy Frontend (Mở Terminal khác)
cd frontend
npm run dev

# Truy cập: http://localhost:3000
# API Docs: http://localhost:8000/docs
# Neo4j Browser: http://localhost:7474
\end{lstlisting}

\subsection{Đánh giá Hệ thống}
\begin{lstlisting}[language=bash, caption=Chạy kiểm thử 50 câu hỏi]
cd backend
.\venv\Scripts\activate
python -X utf8 scripts/evaluate.py
\end{lstlisting}

%% ============================================================
\section{Kết luận}
Dự án \textbf{Graph RAG Chatbot} đã thành công rực rỡ trong việc giải quyết bài toán tư vấn học vụ phức tạp tại ĐHBK Hà Nội. Điểm nhấn mang tính đột phá của dự án là việc khai thác tối đa sức mạnh của thuật toán đồ thị (\textit{Graph Traversal ưu tiên mật độ liên kết}), giúp nâng tỷ lệ tìm trúng đích (Node Hit Rate) lên \textbf{81.7\%} và đạt tỷ lệ trích dẫn đúng Điều khoản tuyệt đối (\textbf{100\%}).

Nhờ sự giám sát chặt chẽ từ System Prompt và trực quan hóa rõ ràng qua giao diện D3.js, người dùng (sinh viên) giờ đây hoàn toàn có thể tin tưởng vào độ chính xác của bot mà không lo sợ rủi ro "ảo giác" (hallucination) thường thấy trên các LLM.

Mặc dù vẫn còn rào cản về phần cứng (đòi hỏi cấu hình máy cá nhân tương đối cao), hệ thống đã chứng minh được tính khả thi xuất sắc của một mô hình Graph RAG chạy hoàn toàn offline (100\% Local, chi phí 0đ) cho các bài toán dữ liệu pháp lý tiếng Việt. Đây là nền tảng vững chắc để tiếp tục mở rộng quy mô tích hợp các tài liệu nội bộ khác trong tương lai.

\vspace{1cm}
\begin{center}
\textit{--- Hết ---}
\end{center}
\end{document}
